% !TeX root = ../../Book5.tex
% 纲要标题：### 16.5 平展等价关系与代数空间
% 纲要位置：几何学计划纲要.md，第 5370 行。
% 正文状态：未撰写；本文件目前是框架。

% 纲要内容开始
% 
% 采用 fppf 层、概形可表对角和概形给出的平展满射图册定义；平展等价关系的层商是本节的构造主线。一般模空间存在定理不由定义自动得到。[R10]
% 
% * 平展等价关系 $R\rightrightarrows U$ 与朴素商的层化
% * 代数空间的层、可表对角与平展满射图册
% * $U\times_XU$ 与图册等价关系的恢复
% * 态射、纤维积及基变换的最低工具
% * 由图册检测局部性质的条件
% * 概形商、代数空间商与非概形模型的区别
% 
% 纲要内容结束

\section{平展等价关系与代数空间}
\label{g5:sec:1605}

待撰写

\subsection{\texorpdfstring{平展等价关系 \(R\rightrightarrows U\) 与朴素商的层化}{平展等价关系 R⇉U 与朴素商的层化}}
\label{g5:subsec:1605:01}
待撰写

\subsection{代数空间的层、可表对角与平展满射图册}
\label{g5:subsec:1605:02}
待撰写

\subsection{\texorpdfstring{\(U\times_XU\) 与图册等价关系的恢复}{U×\_X U 与图册等价关系的恢复}}
\label{g5:subsec:1605:03}
待撰写

\subsection{态射、纤维积及基变换的最低工具}
\label{g5:subsec:1605:04}
待撰写

\subsection{由图册检测局部性质的条件}
\label{g5:subsec:1605:05}
待撰写

\subsection{概形商、代数空间商与非概形模型的区别}
\label{g5:subsec:1605:06}
待撰写
% 纲要条目排为编号小节标题，并进入目录与书签；说明段仍保留于纲要注释中。
